\subsubsection{ATT stability}
\label{sec:hull_att_stability}

Controlled-input (ATT) stability asks whether SMILE attributions remain similar after a
minor, meaning-preserving change to the explanation input.
The perturbation used here expands the explanation reference by one real Lower Super
Output Area (LSOA) just outside Hull in each cardinal direction, one addition at a time
(Table~\ref{tab:hull_att_stability_neighbors}).
Police.uk street counts and the same Census 2021 / OpenStreetMap covariates used for
Hull LSOAs are collected for each neighbour.
Predictive models remain trained on Hull only; only the SMILE reference is expanded.
This quantity is distinct from multi-run \emph{consistency}, which varies explanation
seeds rather than the input.

\begin{table}[htbp]
\centering
\caption{Neighbour LSOAs added one at a time to the SMILE explanation reference.
Codes are Census 2021 LSOA identifiers.
South does not share a Hull boundary; it is the Hessle unit immediately beyond the
LSOA containing Hessle Golf Club.}
\label{tab:hull_att_stability_neighbors}
\begin{tabular}{lll}
\toprule
Direction & LSOA code & Name \\
\midrule
North & \texttt{E01012953} & East Riding of Yorkshire 023A \\
East  & \texttt{E01013094} & East Riding of Yorkshire 029A \\
South & \texttt{E01013009} & East Riding of Yorkshire 034B \\
West  & \texttt{E01013112} & East Riding of Yorkshire 025D \\
\bottomrule
\end{tabular}
\end{table}
\FloatBarrier

The Hull-only baseline is the published two-branch GeoSMILE ranking
(Feature $1024$ masks; Geo $2000$ LSOA perturbations; SMILE seed $42$).
Expansions reuse that design exactly: they do not shrink the mask or perturbation
budget, and they keep the same seed.

Agreement is summarised by
(i)~Jaccard overlap of the top-$k$ attribution set
($k{=}10$ for Feature; $k{=}20$ for Geo) and
(ii)~Spearman correlation of absolute coefficients over the full attribution vector.
The added LSOA is treated as a real Geo player with a fitted coefficient, not forced
to zero.
Hull LSOAs keep the baseline column order and the same perturbation draws, so any
change in the Geo top set would be due to the extra unit rather than a reshuffle of
the random design.

\begin{table}[htbp]
\centering
\caption{Controlled-input (ATT) stability after adding one real LSOA just outside Hull
in each cardinal direction (one-at-a-time; models remain Hull-trained).
Baseline is the published Hull-only Feature top-$10$ and Geo top-$20$ ranking
($1024$ Feature masks; $2000$ Geo perturbations; SMILE seed $42$).
Accepted if Hull-row prediction Spearman $\ge 0.99$.
Top-$k$ Jaccard compares the expanded top set with that published baseline
(Feature $k{=}10$; Geo $k{=}20$).
$|\mathrm{coef}|$ Spearman is the rank correlation of absolute SMILE coefficients
over the full attribution vector.
Pred.\ Spearman is agreement of model outputs on the original Hull test rows.}
\label{tab:hull_att_stability}
\begin{tabular}{llcccc}
\toprule
Addition & Branch & Top-$k$ Jaccard & $|\mathrm{coef}|$ Spearman & Pred.\ Spearman & Accepted \\
\midrule
North & Feature & $1.000$ & $1.000$ & $1.000$ & Yes \\
North & Geo     & $1.000$ & $0.999$ & $1.000$ & Yes \\
East  & Feature & $1.000$ & $0.999$ & $1.000$ & Yes \\
East  & Geo     & $1.000$ & $0.999$ & $1.000$ & Yes \\
South & Feature & $1.000$ & $1.000$ & $1.000$ & Yes \\
South & Geo     & $1.000$ & $1.000$ & $1.000$ & Yes \\
West  & Feature & $1.000$ & $0.999$ & $1.000$ & Yes \\
West  & Geo     & $1.000$ & $0.999$ & $1.000$ & Yes \\
\midrule
Mean  & Feature & $1.000$ & $0.999$ & $1.000$ & $4/4$ \\
Mean  & Geo     & $1.000$ & $0.999$ & $1.000$ & $4/4$ \\
\bottomrule
\end{tabular}
\end{table}

Feature top-$10$ sets are unchanged under this expansion
(Jaccard $1.000$; $|\mathrm{coef}|$ Spearman $0.999$;
Table~\ref{tab:hull_att_stability}).
The same ten predictors remain at the head of the ranking as in the published
global Feature attribution, led by \texttt{crime\_rolling\_mean\_6}
(Figure~\ref{fig:hull_att_stability_feature}).
Geo top-$20$ \emph{membership} is unchanged in every direction (Jaccard $1.000$;
$20/20$ overlap; Table~\ref{tab:hull_att_stability}).
What does change is the order of those same twenty LSOAs: adjacent ranks can swap
when absolute coefficients are close, so map labels $1$--$20$ need not be identical
to the Hull-only baseline even though no new Hull LSOA enters the set and none of
the original twenty leave it.
Figure~\ref{fig:hull_att_stability_geo_base_west} shows this for the west addition,
the extra unit closest to the cutoff: the highlighted top-$20$ set is the same as
the published Hull-only map, while several nearby rank labels move.
The full-vector $|\mathrm{coef}|$ Spearman remains $0.999$, so these are local
re-orderings rather than a change in which places are attributed as important.
The added LSOA never enters the Geo top-$20$
(Table~\ref{tab:hull_att_stability_extra}):
north rank $120$, east $110$, south $112$, and west $46$ of $169$ units.
West is the closest extra unit, still well outside the top-$20$ cutoff.
These are association diagnostics for the fitted explanation, not causal effects.
Figure~\ref{fig:hull_att_stability_geo} shows the four direction maps from the same
runs.

\begin{table}[htbp]
\centering
\caption{Geo rank of the added LSOA after it is included as a real attribution player
with a fitted coefficient.
Rank $1$ is the largest absolute Geo coefficient among $169$ units
($168$ Hull LSOAs plus the neighbour).
None of the four additions enter the Geo top-$20$.}
\label{tab:hull_att_stability_extra}
\begin{tabular}{llccc}
\toprule
Addition & LSOA code & Extra rank & Extra $\lvert\mathrm{coef}\rvert$ & In top-$20$ \\
\midrule
North & \texttt{E01012953} & $120$ & $0.317$ & No \\
East  & \texttt{E01013094} & $110$ & $0.433$ & No \\
South & \texttt{E01013009} & $112$ & $0.396$ & No \\
West  & \texttt{E01013112} & $46$  & $0.945$ & No \\
\bottomrule
\end{tabular}
\end{table}

\begin{figure}[htbp]
\centering
\includegraphics[width=\textwidth]{{Chapter4/Figs/stability runs geo base vs west addition version}.pdf}
\caption{Geo SMILE top-$20$ LSOA importance maps for the published Hull-only
baseline (left) and after adding the west neighbour \texttt{E01013112} (right).
The same twenty Hull LSOAs remain in the top set (Jaccard $=1.000$);
numeric labels $1$--$20$ can differ because nearby ranks swap when
$\lvert\mathrm{coef}\rvert$ values are close.
The added west LSOA is outlined and does not enter the top-$20$ (rank $46$ of
$169$).
Colour encodes $\lvert\mathrm{geo\ coefficient}\rvert$ on the highlighted top-$20$
set.
Published SMILE design ($2000$ Geo perturbations; seed $42$).
Not causal.}
\label{fig:hull_att_stability_geo_base_west}
\end{figure}

\begin{figure}[htbp]
\centering
\includegraphics[width=\textwidth]{{Chapter4/Figs/stability runs addition version FEATURE}.pdf}
\caption{Feature SMILE importance rankings before and after a prediction-preserving
input expansion, using the published SMILE design ($1024$ masks; seed $42$).
Left: published Hull-only global Feature ranking (top-$10$ highlighted).
Right: ranking after adding the west neighbour
(\texttt{E01013112}; top-$10$ Jaccard vs.\ baseline $=1.000$).
Bars show absolute SMILE coefficients; higher bars indicate stronger attribution
magnitude.}
\label{fig:hull_att_stability_feature}
\end{figure}

\begin{figure}[htbp]
\centering
\includegraphics[width=\textwidth, trim=9cm 0cm 5cm 0cm, clip]{{Chapter4/Figs/4 direction addition}.pdf}
\caption{Geo SMILE top-$20$ LSOA importance maps after each one-at-a-time
cardinal-neighbour addition (north, east, south, west), using the published
SMILE design ($2000$ Geo perturbations; seed $42$).
Predictive models remain Hull-trained; only the explanation reference is expanded
by one real LSOA just outside Hull.
Coloured units are the $20$ largest absolute Geo SMILE coefficients; numeric labels
are importance ranks ($1$~=~largest).
The added neighbour is outlined and does not enter the top-$20$ in any direction.
The same twenty Hull LSOAs remain in the top set in all four panels
(Geo top-$20$ Jaccard vs.\ the published Hull-only baseline $=1.000$);
numeric labels may differ because nearby ranks can swap when
$\lvert\mathrm{coef}\rvert$ values are close.
Hull-row prediction Spearman $=1.000$.
These are association diagnostics for the fitted explanation, not causal effects.}
\label{fig:hull_att_stability_geo}
\end{figure}
